\chapter{Análisis curricular}
\label{cap:capitulo3}




\section{Marco normativo: LOMLOE y LOFP}


La integración de la sostenibilidad digital en el sistema educativo español tiene respaldo explícito en la normativa vigente. La \textbf{LOMLOE} (LO 3/2020) fija entre sus fines el "respeto al medio ambiente y al entorno natural" y establece "el desarrollo sostenible" como principio del sistema (Art. 1.l). La \textbf{LOFP} (LO 3/2022) refuerza este mandato al incorporar la sostenibilidad como eje transversal de toda la Formación Profesional, y el \textbf{Real Decreto 659/2023} la consolida como resultado de aprendizaje transversal en la familia profesional de Informática y Comunicaciones.


Los decretos curriculares de referencia para este análisis son el RD 217/2022 (currículo básico ESO), el RD 243/2022 (Bachillerato), y los decretos de título de los ciclos formativos: RD 450/2010 (DAM), RD 686/2010 (DAW), RD 1629/2009 (ASIR) y RD 1691/2007 (SMR).




\section{Análisis de candidatos curriculares}


La sostenibilidad digital es un tema que \textbf{puede y debe trabajarse en todos los niveles educativos}, adaptando la profundidad y las herramientas a la madurez del alumnado. A continuación se describen las posibilidades en cada etapa.


\subsection{Candidatos en Educación Secundaria Obligatoria (ESO)}


\textbf{Marco:} Real Decreto 217/2022, de 29 de marzo.


Las materias de ESO ofrecen el primer contacto posible con la sostenibilidad digital. En \textbf{Tecnología y Digitalización} (1º-2º ESO), el Bloque A "Tecnología y Sociedad" incluye referencias explícitas a los ODS y al impacto ambiental de la tecnología, lo que permite trabajar la huella digital desde la concienciación: qué es, cómo se genera, cómo se puede reducir en el día a día (streaming, almacenamiento en la nube, videollamadas). En \textbf{Tecnología} (3º ESO) se puede profundizar en el ciclo de vida de los dispositivos y en el concepto de eficiencia energética. En 4º ESO, dentro de la Opción Tecnológica, \textbf{Digitalización} permite abordar la huella digital y el consumo energético de los servicios digitales, mientras que \textbf{Tecnología} ofrece contexto para trabajar la eficiencia energética del hardware y los sistemas.


El trabajo en ESO es esencialmente conceptual y actitudinal: el objetivo es activar los componentes cognitivo y afectivo de la actitud pro-sostenibilidad \cite{rosenberg1960cognitive} y sembrar la base sobre la que los niveles superiores construirán la acción técnica.


\subsection{Candidatos en Bachillerato}


\textbf{Marco:} Real Decreto 243/2022, de 5 de abril.


El Bachillerato permite avanzar hacia la \textbf{medición simple y el análisis crítico}. En particular, \textbf{Tecnología e Ingeniería I} (1º Bach) y \textbf{Tecnología e Ingeniería II} (2º Bach) son los contextos más adecuados para integrar prácticas de sostenibilidad digital: al orientarse a las especialidades técnicas, su alumnado ya maneja conceptos de programación, algoritmos y estructuras de datos que permiten introducir la comparación de eficiencia energética y el análisis del impacto ambiental del software. Herramientas como Website Carbon Calculator o Carbonalyser son accesibles sin conocimientos avanzados y permiten cuantificar el impacto de páginas web y hábitos de navegación.


El Bachillerato es un eslabón imprescindible en la cadena formativa: consolida la comprensión conceptual e introduce la medición, preparando el terreno para la acción técnica avanzada que se desarrollará en FP.


\subsection{Candidatos en Formación Profesional — Familia Informática y Comunicaciones}


\textbf{Marco:} RD 450/2010 (DAM), RD 686/2010 (DAW), RD 1629/2009 (ASIR), RD 1691/2007 (SMR).


La FP es el nivel donde es posible alcanzar el \textbf{máximo grado de acción técnica medible}. En \textbf{CFGM SMR}, el foco natural se sitúa en el hardware: ciclo de vida de dispositivos, consumo energético de equipos, RAEE y virtualización como estrategia de eficiencia. En \textbf{CFGS ASIR}, los módulos de Implantación de Sistemas Operativos y, especialmente, Servicios de Red e Internet permiten trabajar la eficiencia energética de la infraestructura y los centros de datos. En \textbf{CFGS DAW}, los módulos de Desarrollo Web en Entorno Servidor y Despliegue de Aplicaciones Web son contextos idóneos para prácticas de Green Software en entornos web y cloud.


\textbf{CFGS DAM — Módulo Programación (MP0485)} y \textbf{CFGS DAW — Desarrollo Web en Entorno Servidor (MP0612)} son los candidatos con mayor potencial para las prácticas de GSE más avanzadas, incluyendo la medición con CodeCarbon, el estándar SCI y el Carbon-Aware Computing.




\section{Asignatura/módulo seleccionado y justificación}


\subsection{Un enfoque transversal con un foco de profundidad}


Antes de exponer la selección, es fundamental precisar el alcance de esta propuesta. La integración de la sostenibilidad digital en la educación informática \textbf{no es un asunto exclusivo de ningún nivel ni ciclo}. Por el contrario, la naturaleza transversal de la crisis climática y el mandato explícito de la LOMLOE y la LOFP implican que \textbf{los principios de Green IT y la concienciación sobre la huella digital deben estar presentes en todos los niveles del sistema educativo}, adaptados a la madurez cognitiva y técnica del alumnado:


\begin{itemize}
\item En \textbf{ESO}, el trabajo es conceptual y de concienciación: los alumnos aprenden qué es la huella digital, cuál es el impacto de sus hábitos cotidianos (streaming, almacenamiento en la nube, videollamadas), y comprenden los ODS desde su realidad diaria.
\item En \textbf{Bachillerato}, se puede avanzar hacia la medición simple y el análisis crítico: herramientas como Website Carbon Calculator o Carbonalyser permiten cuantificar el impacto de webs y hábitos de navegación sin requerir conocimientos avanzados de programación.
\item En \textbf{FP de Grado Medio (SMR)}, el foco se sitúa en el hardware: ciclo de vida de dispositivos, consumo energético de equipos, RAEE, y virtualización como estrategia de eficiencia.
\item En \textbf{FP de Grado Superior (DAM, DAW, ASIR)}, es posible alcanzar el nivel más avanzado: la \textbf{acción técnica medible}, donde el alumnado no solo conoce y mide el problema sino que lo \textbf{resuelve} aplicando principios de Ingeniería de Software Verde con herramientas profesionales reales.
\end{itemize}


La presente propuesta didáctica \textbf{se centra en el nivel de Grado Superior de FP}, no porque los demás niveles sean menos importantes —son igualmente necesarios en una estrategia educativa coherente con la transversalidad de la EDS—, sino porque este contexto ofrece las condiciones óptimas para demostrar la viabilidad y el impacto real de las prácticas de GSE: alumnado con formación técnica suficiente, herramientas de medición directamente aplicables, y un horizonte profesional inmediato que convierte el aprendizaje en acción transformadora real. Las prácticas diseñadas en este TFM están concebidas, además, para ser \textbf{adaptadas y escaladas} a otros niveles educativos, constituyendo un marco metodológico transferible que el docente puede simplificar o complejizar según el contexto.


\subsection{Selección: CFGS Desarrollo de Aplicaciones Multiplataforma (DAM) — Módulo Profesional MP0485: Programación — 1º curso}


Se selecciona \textbf{DAM — Programación} frente a DAW por tres razones principales. Primero, es el módulo nuclear de 1º (256 h), lo que permite introducir los principios de GSE desde el inicio de la formación técnica del alumnado. Segundo, trabaja con Python/Java, lenguajes para los que existen herramientas de medición maduras (CodeCarbon, JoularJX); en DAW el stack web añade complejidad de medición que dificulta las prácticas de aula. Tercero, el alumnado de DAM se incorpora al mercado laboral como desarrollador en 12-18 meses, lo que convierte el aprendizaje en impacto ambiental directo y cuantificable.


\begin{table}[htbp]
\centering
\begin{tabular}{|l|p{10cm}|}
\hline
\textbf{Campo} & \textbf{Valor} \\
\hline
Ciclo & CFGS Desarrollo de Aplicaciones Multiplataforma (DAM) \\
\hline
Módulo & Programación (MP0485) — 1º curso \\
\hline
Carga horaria & 256 h (\textasciitilde{}8 h/semana) \\
\hline
Marco normativo & Real Decreto 450/2010, de 16 de abril \\
\hline
\end{tabular}
\caption{Selección: CFGS Desarrollo de Aplicaciones Multiplataforma (DAM) — Módulo Profesional MP0485: Programación — 1º curso}
\label{tab:seleccion-cfgs-desarrollo-de-aplicaciones-multipla}
\end{table}




\section{Resultados de aprendizaje y puntos de inserción}


El RD 450/2010 estructura MP0485 en \textbf{nueve Resultados de Aprendizaje (RA)}:


\begin{table}[htbp]
\centering
\begin{tabular}{|l|p{10cm}|}
\hline
\textbf{RA} & \textbf{Descripción} \\
\hline
RA1 & Reconoce la estructura de un programa informático, identificando y relacionando los elementos propios del lenguaje de programación utilizado. \\
\hline
RA2 & Escribe y prueba programas sencillos, reconociendo y aplicando los fundamentos de la programación orientada a objetos. \\
\hline
RA3 & Escribe y depura código, analizando y utilizando las estructuras de control del flujo. \\
\hline
RA4 & Desarrolla programas organizados en clases analizando y aplicando los principios de la programación orientada a objetos. \\
\hline
RA5 & Realiza operaciones de entrada y salida de información, utilizando procedimientos específicos del lenguaje y librerías de clases. \\
\hline
RA6 & Escribe programas que manipulan información seleccionando y utilizando tipos avanzados de datos. \\
\hline
RA7 & Desarrolla programas aplicando características avanzadas de los lenguajes orientados a objetos y del entorno de programación. \\
\hline
RA8 & Utiliza bases de datos orientadas a objetos, analizando sus características y aplicando técnicas para gestionar la persistencia de los objetos. \\
\hline
RA9 & Gestiona información almacenada en bases de datos relacionales manteniendo la integridad y consistencia de los datos. \\
\hline
\end{tabular}
\caption{Resultados de aprendizaje y puntos de inserción}
\label{tab:resultados-de-aprendizaje-y-puntos-de-insercion}
\end{table}


\subsection{Resultados de aprendizaje seleccionados para las prácticas Green IT}


Las cuatro prácticas de esta propuesta se insertan en los RA con mayor alineación técnica con los principios GSE:


\begin{table}[htbp]
\centering
\begin{tabularx}{\textwidth}{|X|X|X|X|}
\hline
\textbf{Práctica} & \textbf{RA} & \textbf{Contenido Green IT} & \textbf{Principio GSE} \\
\hline
P1 — Benchmark energético & \textbf{RA3} (estructuras de control) & Comparación del consumo de CPU entre algoritmos equivalentes & Principio I — Eficiencia Energética \\
\hline
P2 — Memoria y objetos & \textbf{RA4} (POO) & Instanciación de objetos y huella de memoria & Principio III — Hardware eficiente \\
\hline
P3 — E/S eficiente & \textbf{RA5} (entrada/salida) & Streaming vs. carga completa de datos & Principio I — Minimización de transferencia \\
\hline
P4 — Carbon-Aware & \textbf{RA9} (estructuras dinámicas) & Cola de tareas con conciencia de intensidad de carbono & Principio II — Carbon-Aware Computing \\
\hline
\end{tabularx}
\caption{Resultados de aprendizaje seleccionados para las prácticas Green IT}
\label{tab:resultados-de-aprendizaje-seleccionados-para-las-p}
\end{table}


RA3 y RA5 se trabajan en el primer trimestre, cuando el alumnado está construyendo sus primeros hábitos de codificación; RA4 en el segundo, al profundizar en POO; y RA9 al final del curso, cuando ya domina las estructuras de datos y puede integrar APIs externas.




\section{Competencias a desarrollar}


La propuesta didáctica trabaja de forma integrada competencias de tres marcos de referencia:


\subsection{Competencias clave LOMLOE (Real Decreto 217/2022, Anexo I)}


\begin{table}[htbp]
\centering
\begin{tabular}{|l|p{10cm}|}
\hline
\textbf{Competencia Clave} & \textbf{Cómo se trabaja} \\
\hline
\textbf{Competencia Digital (CD)} — Área 4: Seguridad & Protección del entorno digital incluyendo impacto ambiental; gestión responsable del consumo energético. \\
\hline
\textbf{Competencia en Ciencia, Tecnología e Ingeniería (STEM)} & Resolución de problemas técnicos reales con criterio de eficiencia energética como variable de calidad. \\
\hline
\textbf{Competencia Ciudadana (CC)} & Comprensión de la responsabilidad del desarrollador de software en el contexto de la crisis climática global. \\
\hline
\textbf{Competencia Personal, Social y de Aprender a Aprender (CPSAA)} & Reflexión metacognitiva sobre los hábitos de programación y su impacto medioambiental. \\
\hline
\end{tabular}
\caption{Competencias clave LOMLOE (Real Decreto 217/2022, Anexo I)}
\label{tab:competencias-clave-lomloe-real-decreto-217-2022-an}
\end{table}


\subsection{GreenComp — Marco Europeo de Competencia en Sostenibilidad (2022)}


\begin{table}[htbp]
\centering
\begin{tabularx}{\textwidth}{|X|X|X|}
\hline
\textbf{Área GreenComp} & \textbf{Competencia} & \textbf{Cómo se trabaja} \\
\hline
\textbf{Conceptualización de un futuro verde} & \textit{Literacy in sustainability} & Comprensión del impacto ambiental de las TIC y del rol del software. \\
\hline
\textbf{Valoración de la sostenibilidad} & \textit{Embodying sustainability values} & Incorporar la eficiencia energética como criterio de calidad del código, no como añadido opcional. \\
\hline
\textbf{Acción para la sostenibilidad} & \textit{Taking action} & Implementación real de código con menor huella de carbono medida con SCI/CodeCarbon. \\
\hline
\textbf{Acción para la sostenibilidad} & \textit{Individual initiative} & El alumnado propone mejoras de eficiencia de forma autónoma en su propio código. \\
\hline
\end{tabularx}
\caption{GreenComp — Marco Europeo de Competencia en Sostenibilidad (2022)}
\label{tab:greencomp-marco-europeo-de-competencia-en-sostenib}
\end{table}


\subsection{DigComp 2.2 (2022) — Competencia Digital}


\begin{table}[htbp]
\centering
\begin{tabularx}{\textwidth}{|X|X|X|}
\hline
\textbf{Área DigComp} & \textbf{Competencia} & \textbf{Relación con Green IT} \\
\hline
Área 4: Seguridad & 4.4 \textit{Protecting the environment} & Uso responsable de los dispositivos; impacto ambiental de las prácticas digitales. \\
\hline
Área 5: Resolución de problemas & 5.2 \textit{Identifying needs and technological responses} & Identificar soluciones técnicas (GSE) a problemas ambientales reales. \\
\hline
\end{tabularx}
\caption{DigComp 2.2 (2022) — Competencia Digital}
\label{tab:digcomp-2-2-2022-competencia-digital}
\end{table}




\section{Síntesis: justificación curricular completa}


La selección del módulo \textbf{Programación (MP0485) del CFGS DAM} como contexto principal de esta propuesta didáctica responde a una convergencia de factores normativos, técnicos, pedagógicos y estratégicos.


En primer lugar, cabe reiterar que \textbf{la sostenibilidad digital es un imperativo transversal} que debe impregnar toda la cadena educativa del sistema de informática: desde la ESO, donde se siembran actitudes y conciencia, pasando por el Bachillerato, donde se desarrolla el análisis crítico, hasta la FP, donde se consolida la competencia técnica de acción. Este TFM no propone sustituir ni ignorar esos otros niveles: propone un modelo de máxima profundidad que puede servir de referencia y del que se pueden derivar versiones adaptadas a cada etapa.


Dicho esto, la elección de DAM — Programación como foco del trabajo se justifica por cuatro razones articuladas:


\begin{enumerate}
\item \textbf{Normativo:} el RD 659/2023 y la LOFP establecen explícitamente la sostenibilidad como
\end{enumerate}

resultado de aprendizaje transversal en la familia de Informática y Comunicaciones, dotando de cobertura legal y mandato institucional a la integración propuesta.


\begin{enumerate}
\item \textbf{Técnico:} el módulo de Programación proporciona el lenguaje (Python/Java), las
\end{enumerate}

herramientas (CodeCarbon, Electricity Maps API, SCI) y los conceptos (algoritmos, estructuras de datos, E/S) necesarios para trabajar los Principios I y II de la GSE con el rigor que requiere una propuesta didáctica de máxima complejidad. Ningún otro nivel del sistema educativo ofrece este grado de acceso simultaneo a herramientas profesionales reales y tiempo curricular suficiente para usarlas con profundidad.


\begin{enumerate}
\item \textbf{Pedagógico:} la inserción en 1º de DAM garantiza que los principios de Green IT se
\end{enumerate}

conviertan en parte de la identidad profesional del alumnado desde el inicio de su formación, con capacidad de transferencia a todos los módulos posteriores del ciclo y a los proyectos de fin de ciclo. No se trata de añadir un tema sobre sostenibilidad al final del curso: se trata de hacer del criterio de eficiencia energética una \textbf{variable de calidad del código desde el primer ejercicio}.


\begin{enumerate}
\item \textbf{Estratégico e impacto inmediato:} el alumnado de 1º DAM se incorporará al mercado
\end{enumerate}

laboral como desarrollador de software en un horizonte de 12-18 meses. A diferencia de otros niveles educativos donde el impacto social del aprendizaje es diferido en el tiempo, \textbf{formar a futuros desarrolladores en GSE produce un retorno ambiental directo y a corto plazo}: el código que escriben en producción consume (o no) energía real. Un desarrollador que ha interiorizado el SCI, que ha medido con CodeCarbon y que conoce el Carbon-Aware Computing no toma las mismas decisiones de diseño que uno que no ha recibido esta formación. La escala del sector —millones de desarrolladores generando software que se ejecuta en millones de dispositivos— convierte esta formación en una de las intervenciones educativas con mayor potencial de impacto ambiental real por hora de docencia invertida.
